\section{Coefficients from the complete theta source}\label{sec:coefficients}

The preceding identities isolate what must be estimated. We now
express the coefficients using the actual source and its primitive,
retaining all signed terms and endpoint contributions.

\subsection{Signed source moments}
At \(\lambda_* =\lambda_0-\tau_*\), define
\begin{equation}\label{eq:coefficient-moments}
 C_j=\int_0^\infty u^j F_{\lambda_*}(u)\cos(x_*u)\dd u,
 \qquad
 S_j=\int_0^\infty u^j F_{\lambda_*}(u)\sin(x_*u)\dd u.
\end{equation}
The complete bounds of Section~\ref{sec:source} justify
differentiation under each integral. At a collision they give
\begin{align}
 C_0&=S_1=0,\qquad a=-C_2,\quad b=S_3,\quad c=C_4,
 \label{eq:moment-jets}\\
 \beta&=-\frac{S_3}{C_2},\qquad
 \gamma=-\frac{C_2C_4+S_3^2}{C_2^2}.
 \label{eq:moment-ratios}
\end{align}
These formulas hold for the complete arithmetic source. Positivity
of \(F_{\lambda_*}\) does not make the signed cosine moments a
positive moment matrix. Separate absolute bounds on \(C_2,S_3,C_4\)
do not bound their ratios without a lower bound on \(|C_2|\).
A coupled estimate or an independently controlled exterior factor
would provide additional information; the formulas alone do not.

\subsection{The half-line primitive and its endpoint}
Use the restriction to \(u\ge0\) of the function \(A\) in
Section~\ref{sec:source}, now denoted by
\begin{equation}\label{eq:p-primitive}
 p(u)=\ee^u\sum_{n\ge1}\ee^{-\pi n^2\ee^{4u}},\qquad
 p''-p=8\Phi,\qquad p'(0)=-\tfrac12.
\end{equation}
The differential identity follows by termwise differentiation;
differentiating \(A(-u)-A(u)=\sinh u\) at zero proves the endpoint.
The same double-exponential majorants as before, with the corresponding
lower-degree polynomials, justify all integrations below. Set
\begin{equation}\label{eq:J-primitive}
 J_\lambda(x)=\int_0^\infty p(u)\ee^{\lambda u^2}\cos(xu)\dd u.
\end{equation}
Here \(J_\lambda\) is distinct from the absolute majorants
\(J_{r,d}\) in \eqref{eq:J}.

\begin{proposition}[Exact primitive transfer]\label{prop:primitive-transfer}
For every finite real \(\lambda,x\),
\begin{equation}\label{eq:primitive-transfer}
 8H_\lambda(x)=\tfrac12+(2\lambda-1-x^2)J_\lambda(x)
       +4\lambda xJ_\lambda'(x)-4\lambda^2J_\lambda''(x).
\end{equation}
For \(n\ge2\), with primes denoting spatial derivatives,
\begin{align}
 8H_\lambda^{(n)}={}&-4\lambda^2J_\lambda^{(n+2)}
 +4\lambda xJ_\lambda^{(n+1)}
 +(2\lambda-1-x^2+4\lambda n)J_\lambda^{(n)}\notag\\
 &-2nxJ_\lambda^{(n-1)}-n(n-1)J_\lambda^{(n-2)}.
 \label{eq:primitive-jets}
\end{align}
In particular \(a,b,c\) require primitive derivatives only through
orders four, five and six respectively.
\end{proposition}
\begin{proof}
For the smooth decaying tests used here, two integrations by parts give
\begin{equation}\label{eq:primitive-IBP}
 8\int_0^\infty\Phi f\dd u
 =\tfrac12 f(0)+p(0)f'(0)
       +\int_0^\infty p(f''-f)\dd u.
\end{equation}
Indeed the endpoint term is
\([p'f-pf']_0^\infty=-p'(0)f(0)+p(0)f'(0)\).
For \(f(u)=\ee^{\lambda u^2}\cos(xu)\), \(f(0)=1\), \(f'(0)=0\),
and
\[
 f''-f=\ee^{\lambda u^2}
 \bigl((2\lambda-1-x^2+4\lambda^2u^2)\cos(xu)
                   -4\lambda xu\sin(xu)\bigr).
\]
Replacing the last two integrals by \(-J_\lambda''\) and
\(-J_\lambda'\) proves \eqref{eq:primitive-transfer}.
Leibniz' rule for the quadratic and linear coefficients proves
\eqref{eq:primitive-jets}. All differentiations and vanishing
boundary terms are uniform on compact parameter sets by the paid
primitive tails.
\end{proof}

At \(\lambda=0\) this reduces to
\(8H_0=\tfrac12-(1+x^2)J_0\). A hypothetical collision on that
slice would consequently satisfy
\[
 J_0(x_*)=\frac{1}{2(1+x_*^2)},\qquad
 J_0'(x_*)=-\frac{x_*}{(1+x_*^2)^2}.
\]
The higher derivatives still contain the coefficient information.
The constant \(1/2\) in \eqref{eq:primitive-transfer} cannot be
dropped. Nor does positivity of \(p\) assert positivity of its
oscillatory transform or of the displayed signed combinations.

This half-line primitive is different from the smooth full-line
Green field \(G_\lambda\) in \eqref{eq:G-transform}.
In general \(G_\lambda\ne\ee^{\lambda u^2}G_0\); the latter is
not even integrable when \(\lambda>0\), by \eqref{eq:G-tail}.
The primitive identity is an exact alternative source representation,
with no estimating advantage claimed merely from the transfer.

\subsection{Recovering the exterior coefficients on a contour}
The local factor also has an interface that avoids division by the
vanishing value \(h(\tau_*,x_*)\). Let \(\Gamma\) be a positively
oriented simple contour bounding a disk that contains exactly the
two-root cluster, with multiplicity, and no boundary zero. Factor
\(H=P U\) there, with \(P\) the monic quadratic and \(U\) nonzero.
For \(z\) inside the contour define
\begin{equation}\label{eq:contour-A}
 \mathcal A_H(z)=\frac{1}{2\pi\ii}\oint_\Gamma
       \frac{H'(\zeta)/H(\zeta)}{\zeta-z}\dd\zeta.
\end{equation}
Then
\begin{equation}\label{eq:contour-exterior}
 \mathcal A_H(z)=\frac{U'(z)}{U(z)},\qquad
 \ell=\mathcal A_H(x_*),\quad \kappa=\mathcal A_H'(x_*).
\end{equation}
To prove this, subtract \(P'/P\) from \(H'/H\) and apply
Cauchy's formula to \(U'/U\). For \(z\) off the roots, each
term \(1/(\zeta-r)\) in \(P'/P\) contributes two opposite
residues to \eqref{eq:contour-A}, at \(r\) and \(z\).
Its integral is zero. Analytic continuation gives the formula at
the roots as well. Differentiating under the contour integral
gives the denominator \((\zeta-z)^2\) for \(\kappa\).
This is the existing local exterior-force interface, reused here
for the critical coefficients \cite{RTLocalTransfer}.

For clarity, its error payment can be stated explicitly. Suppose
on \(\Gamma\), \(|H|\ge m>0\),
\(|\widehat H-H|\le e_0<m\), \(|\widehat H'-H'|\le e_1\),
and \(|H'|\le M_1\). Then
\[
 \left|\frac{\widehat H'}{\widehat H}-\frac{H'}H\right|
 \le \frac{e_1}{m-e_0}+\frac{M_1e_0}{m(m-e_0)}=:D.
\]
If the length of \(\Gamma\) is \(L_\Gamma\) and its distance
from \(z\) is at least \(d>0\), the errors in \(\mathcal A_H(z)\)
and \(\mathcal A_H'(z)\) are at most
\(L_\Gamma D/(2\pi d)\) and \(L_\Gamma D/(2\pi d^2)\).
These follow directly by integrating the ratio error. The method
needs a boundary margin, a paid cluster count and complete source
errors; it does not need an internal pair-gap lower bound.
Its advantage, if any, is a quantitative comparison on a stated
region, not a consequence of introducing the contour.
